% EDITORIAL QUESTION: Quais definições operacionais padronizam a comunicação entre setters, gestores e praticantes?
\section{Glossário Técnico e Termos Operacionais}

\begin{description}
    \item[Beta / Beta Break:] Sequência de movimentos planejada pelo \textit{setter} para resolver um bloco; \textit{beta break} ocorre quando o escalador encontra solução alternativa não prevista que reduz a dificuldade pretendida.
    \item[Conquistador / Abertura de Via:] Termos reservados neste documento exclusivamente para a escalada em rocha natural (conforme ética da CBME DT-2016/10).
    \item[Crux:] Lance ou sequência de movimentos mais difícil de uma via ou boulder.
    \item[Forerunning:] Processo coletivo de testagem prévia de blocos recém-montados por múltiplos \textit{setters} de estaturas distintas antes da liberação ao público.
    \item[Grau Trabalhado (``Decotar''):] Execução após ensaio prévio dos movimentos (\textit{redpoint}), em contraste com a escalada \textit{à vista} (\textit{onsight}) exigida como referência na CBME DT-2016/01.
    \item[Setter (ou Routesetter):] Profissional ou membro técnico responsável por projetar, montar e calibrar percursos em muros artificiais.
    \item[Start Tag:] Etiqueta fixada na agarra de largada contendo circuito, cor, metadados e subgrau.
\end{description}
